\documentclass[aspectratio=169]{beamer}

\usetheme{Madrid}
\usecolortheme{seahorse}

\usepackage{graphicx}
\usepackage{booktabs}
\usepackage{amsmath}
\usepackage{amssymb}
\usepackage{tabularx}
\usepackage{array}
\usepackage{ragged2e}
\usepackage{tikz}
\usepackage{pgfplots}
\pgfplotsset{compat=1.18}

\graphicspath{{./}{figures/}}

% The Chapter 15 RMarkdown notes generate the figure files used below.
% A labelled placeholder is shown if the slides are compiled before the
% notes have been knitted.
\newcommand{\lectureimage}[2][0.90\linewidth]{%
  \IfFileExists{#2}{%
    \includegraphics[width=#1,height=0.74\textheight,keepaspectratio]{#2}%
  }{%
    \fbox{%
      \parbox[c][0.38\textheight][c]{0.88\linewidth}{%
        \centering\small
        Figure file: \texttt{\detokenize{#2}}\\[0.15cm]
        Generate this figure by knitting the Chapter 15 RMarkdown notes.
      }%
    }%
  }%
}

\title[Chapter 15: Inference in Practice]{\textbf{Chapter 15: Inference in Practice}}
\subtitle{The Practice of Statistics in the Life Sciences \\ by Brigitte Baldi and David S. Moore}
\author{Instructor: Ali Raisolsadat}
\date{Part III: Foundations of Inference}

\setbeamertemplate{navigation symbols}{}
\setbeamertemplate{itemize items}[circle]
\setbeamertemplate{enumerate items}[default]

\begin{document}

% =============================================================
% INTRODUCTION
% =============================================================

% SLIDE 1
\begin{frame}
  \titlepage
\end{frame}

% SLIDE 2
\begin{frame}{Why Inference Requires Judgment}
Chapter 14 introduced confidence intervals and hypothesis tests under simple conditions.

\vspace{0.25cm}
The formulas are only one part of a statistical analysis. Reliable conclusions also depend on how the data were collected, whether the model is reasonable and whether the numerical result matters in context.

\vspace{0.25cm}
\begin{block}{Central Idea}
Statistical methods are useful when their assumptions fit the study and their conclusions are interpreted with appropriate caution.
\end{block}

\vspace{0.15cm}
This chapter develops the practical judgment needed to use inference responsibly.
\end{frame}

% SLIDE 3
\begin{frame}{Chapter Roadmap}
The chapter focuses on four practical questions.

\begin{enumerate}
  \item When are confidence intervals and hypothesis tests trustworthy?
  \item How do confidence level, variability and sample size affect a confidence interval?
  \item What do statistical significance and a \textit{P}-value establish, and what do they leave unanswered?
  \item How can sample size be planned before data collection begins?
\end{enumerate}

\vspace{0.25cm}
The first part of the chapter reviews the general theory. The second part applies these ideas to examples from science and public research.
\end{frame}

% =============================================================
% THEORY: CONDITIONS FOR INFERENCE
% =============================================================

\section{Conditions for Inference}

% SLIDE 4
\begin{frame}{The Simple Conditions for One-Sample $z$ Procedures}
The procedures introduced in Chapter 14 assumed that

\begin{itemize}
  \item The observations came from a simple random sample
  \item The population distribution was Normal, or the sampling distribution of $\bar{x}$ was approximately Normal
  \item The population standard deviation $\sigma$ was known
\end{itemize}

\vspace{0.25cm}
Under these conditions, the one-sample $z$ procedures use

\[
\bar{x}\pm z^*\frac{\sigma}{\sqrt{n}}
\qquad\text{and}\qquad
z=\frac{\bar{x}-\mu_0}{\sigma/\sqrt{n}}
\]

The assumption that $\sigma$ is known is uncommon in practice. Later chapters replace it with an estimate from the sample.
\end{frame}

% SLIDE 5
\begin{frame}{Where the Data Come From}
Statistical inference relies on probability to describe how a statistic varies across possible samples or treatment assignments.

\vspace{0.2cm}
This reasoning is most reliable when chance enters the study design through

\begin{itemize}
  \item \textbf{Random sampling}, which supports generalization from the sample to the population
  \item \textbf{Random assignment}, which helps balance lurking variables and supports causal conclusions
\end{itemize}

\vspace{0.25cm}
\begin{block}{Practical Question}
Before calculating an interval or a \textit{P}-value, ask how the observations entered the study and which population or treatment comparison the design can support.
\end{block}
\end{frame}

% SLIDE 6
\begin{frame}{Random Sampling and Random Assignment}
The two forms of randomization support different conclusions.

\begin{center}
\begin{tabularx}{0.94\textwidth}{>{\RaggedRight\arraybackslash}p{0.22\textwidth} X X}
\toprule
\textbf{Design feature} & \textbf{Main purpose} & \textbf{Conclusion supported} \\
\midrule
Random sampling & Select a representative sample from a defined population & Generalization to the sampled population \\
Random assignment & Assign experimental units to treatments by chance & Comparison of treatment effects and stronger causal interpretation \\
\bottomrule
\end{tabularx}
\end{center}

\vspace{0.25cm}
A study may use one, both or neither form of randomization. The scope of the conclusion should match the design actually used.
\end{frame}

% SLIDE 7
\begin{frame}{Practical Problems After a Good Design}
Even a carefully designed study can encounter difficulties after data collection begins.

\begin{itemize}
  \item \textbf{Nonresponse} may make respondents systematically different from people who do not respond
  \item \textbf{Attrition} may create bias when participants who leave a study differ from those who remain
  \item \textbf{Measurement problems} may introduce systematic error that the sampling model does not describe
  \item \textbf{Complex sampling} may require methods that account for clustering, stratification or unequal selection probabilities
\end{itemize}

\vspace{0.2cm}
A standard formula cannot correct a fundamental flaw in the study design or a mismatch between the design and the analysis.
\end{frame}

% SLIDE 8
\begin{frame}{The Population Model}
Parametric inference begins with a probability model for the population.

\vspace{0.2cm}
For one-sample inference about a mean, the model describes the distribution of the quantitative variable and identifies the unknown parameter $\mu$.

\vspace{0.2cm}
The Normal condition matters because the procedure uses the sampling distribution of $\bar{x}$. The observations themselves do not need to form a perfect Normal distribution whenever the sample size is large enough for the Central Limit Theorem to apply.

\begin{block}{Distinction}
The shape of the population distribution and the shape of the sampling distribution of $\bar{x}$ are different ideas.
\end{block}
\end{frame}

% SLIDE 9
\begin{frame}{Sample Size and the Normal Approximation}
The amount of non-Normality that a procedure can tolerate depends partly on the sample size.

\begin{itemize}
  \item For a small sample, strong skewness or outliers can make Normal-based inference unreliable
  \item For a moderately sized sample, a roughly symmetric population is often sufficient
  \item For a large sample, the Central Limit Theorem makes the sampling distribution of $\bar{x}$ approximately Normal under broad conditions
\end{itemize}

\vspace{0.2cm}
There is no universal sample-size cutoff that works for every population. Stronger skewness and heavier tails generally require more observations.
\end{frame}

% SLIDE 10
\begin{frame}{Outliers and Nonresistant Statistics}
The sample mean $\bar{x}$ and sample standard deviation $s$ are sensitive to extreme observations.

\vspace{0.2cm}
An outlier can change the estimate, increase the standard error and alter the conclusion of a hypothesis test.

\vspace{0.25cm}
Before performing inference:

\begin{itemize}
  \item Check for data-entry and measurement errors
  \item Examine a dotplot, histogram or another suitable graph
  \item Investigate scientifically plausible reasons for unusual observations
  \item Use an alternative method when a small sample contains severe skewness or influential outliers
\end{itemize}

An observation should not be removed only because its removal produces a preferred result.
\end{frame}

% SLIDE 11
\begin{frame}{A Practical Conditions Checklist}
Before reporting an inferential result, consider the following questions.

\begin{enumerate}
  \item What population or treatment comparison does the study address?
  \item How were individuals selected and, if applicable, assigned to treatments?
  \item Are observations sufficiently independent for the procedure?
  \item Does the distributional model fit the data and sample size?
  \item Are outliers, missing observations or measurement problems present?
  \item Does the chosen procedure match the sampling or experimental design?
\end{enumerate}

\vspace{0.2cm}
These questions determine whether the reported standard error, confidence interval and \textit{P}-value have a defensible interpretation.
\end{frame}

% =============================================================
% THEORY: CONFIDENCE INTERVALS
% =============================================================

\section{How Confidence Intervals Behave}

% SLIDE 12
\begin{frame}{The General Form of a Confidence Interval}
Many confidence intervals have the form

\[
\text{estimate}\pm\text{margin of error}
\]

The margin of error combines two quantities:

\[
\text{margin of error}
=
\text{critical value}\times\text{standard error}
\]

\vspace{0.2cm}
For a one-sample $z$ confidence interval for a population mean,

\[
m=z^*\frac{\sigma}{\sqrt{n}}
\]

The estimate determines the centre of the interval. The margin of error determines its half-width.
\end{frame}

% SLIDE 13
\begin{frame}{Confidence and Precision}
Two desirable features of an interval compete with one another.

\begin{itemize}
  \item A \textbf{high confidence level} gives a procedure with a high long-run coverage rate
  \item A \textbf{small margin of error} gives a more precise estimate of the parameter
\end{itemize}

\vspace{0.25cm}
With the same data, a higher confidence level requires a larger critical value $z^*$ and therefore a wider interval.

\begin{center}
\begin{tabular}{ccc}
\toprule
\textbf{Confidence level} & \textbf{$z^*$} & \textbf{Relative width} \\
\midrule
$90\%$ & $1.645$ & Narrower \\
$95\%$ & $1.960$ & Intermediate \\
$99\%$ & $2.576$ & Wider \\
\bottomrule
\end{tabular}
\end{center}
\end{frame}

% SLIDE 14
\begin{frame}{Population Variability and Precision}
The margin of error increases with the population standard deviation.

\[
m=z^*\frac{\sigma}{\sqrt{n}}
\]

\begin{itemize}
  \item A smaller $\sigma$ means that individual values are more tightly concentrated around the population mean
  \item A larger $\sigma$ produces greater sample-to-sample variation in $\bar{x}$
\end{itemize}

\vspace{0.2cm}
Natural biological variation usually cannot be removed. Measurement error can sometimes be reduced by improving instruments, procedures or training.

\vspace{0.2cm}
The value of $\sigma$ should describe the population and measurement process relevant to the planned study.
\end{frame}

% SLIDE 15
\begin{frame}{Sample Size and the Square-Root Relationship}
For a fixed confidence level and population standard deviation,

\[
m\propto\frac{1}{\sqrt{n}}
\]

Increasing the sample size improves precision, but the improvement follows a square-root relationship.

\begin{center}
\begin{tabular}{ccc}
\toprule
\textbf{Change in sample size} & \textbf{New standard error} & \textbf{Effect on margin} \\
\midrule
$n\rightarrow4n$ & One-half as large & Margin is halved \\
$n\rightarrow9n$ & One-third as large & Margin is divided by $3$ \\
$n\rightarrow100n$ & One-tenth as large & Margin is divided by $10$ \\
\bottomrule
\end{tabular}
\end{center}

Doubling the sample size reduces the margin by a factor of $1/\sqrt{2}\approx0.707$, not by one-half.
\end{frame}

\begin{frame}{Margin of Error and Sample Size}
\centering
\lectureimage[0.88\linewidth]{fig_ch15_margin_sample_size.png}
\end{frame}

% SLIDE 16
\begin{frame}{What the Margin of Error Measures}
The margin of error describes uncertainty from random sampling under the assumed model.

\vspace{0.2cm}
It reflects how far the sample statistic may fall from the population parameter because a different random sample would give a different statistic.

\begin{block}{Covered by the Margin of Error}
Random sample-to-sample variation generated by the stated sampling process
\end{block}

\vspace{0.15cm}
A correctly calculated margin of error can still accompany a misleading conclusion if the observations are biased or the statistical model does not fit the study.
\end{frame}

% SLIDE 17
\begin{frame}{Errors Outside the Margin of Error}
The usual margin of error does not measure uncertainty caused by

\begin{itemize}
  \item Undercoverage or voluntary response
  \item Nonresponse and systematic attrition
  \item Leading questions or inaccurate self-reports
  \item Poorly calibrated instruments or inconsistent measurements
  \item Confounding in an observational study
  \item A statistical model that does not match the design or population
\end{itemize}

\vspace{0.2cm}
These sources of error can be larger than the reported sampling error. A narrow confidence interval does not compensate for biased data.
\end{frame}

% =============================================================
% THEORY: HYPOTHESIS TESTS
% =============================================================

\section{How Hypothesis Tests Behave}

% SLIDE 18
\begin{frame}{Evidence on a Continuum}
A \textit{P}-value measures how unusual the observed statistic would be under the null model, in the direction specified by the alternative hypothesis.

\vspace{0.2cm}
Smaller \textit{P}-values provide stronger evidence against $H_0$, but there is no universal boundary separating meaningful evidence from unimportant evidence.

\begin{block}{Interpretive Caution}
The values $P=0.048$ and $P=0.052$ represent nearly the same strength of evidence, even though a fixed $5\%$ decision rule places them on opposite sides of a threshold.
\end{block}

The exact \textit{P}-value, study design and consequences of a decision should all inform the conclusion.
\end{frame}

% SLIDE 19
\begin{frame}{Type I and Type II Errors}
A hypothesis-test decision can be correct or incorrect.

\begin{center}
\begin{tabularx}{0.92\textwidth}{>{\RaggedRight\arraybackslash}p{0.24\textwidth} X X}
\toprule
 & \textbf{$H_0$ is true} & \textbf{A specified alternative is true} \\
\midrule
Fail to reject $H_0$ & Correct decision & Type II error \\
Reject $H_0$ & Type I error & Correct decision \\
\bottomrule
\end{tabularx}
\end{center}

\[
P(\text{Type I error})=\alpha
\]

\[
P(\text{Type II error at a specified alternative})=\beta
\]

The relative seriousness of the two errors depends on the scientific and practical consequences.
\end{frame}

% SLIDE 20
\begin{frame}{The Alternative Hypothesis Determines the Tail}
The alternative hypothesis must reflect the scientific question established before examining the sample result.

\begin{itemize}
  \item $H_a:\mu>\mu_0$ uses the upper tail
  \item $H_a:\mu<\mu_0$ uses the lower tail
  \item $H_a:\mu\ne\mu_0$ uses both tails
\end{itemize}

\vspace{0.2cm}
For a symmetric test statistic, a one-sided \textit{P}-value is often one-half of the corresponding two-sided \textit{P}-value when the observed effect is in the stated direction.

\vspace{0.2cm}
Choosing the direction after seeing the data understates uncertainty and invalidates the intended error rate.
\end{frame}

% SLIDE 21
\begin{frame}{Statistical Significance and Sample Size}
For a one-sample $z$ test,

\[
z
=
\frac{\bar{x}-\mu_0}{\sigma/\sqrt{n}}
=
\frac{\text{observed effect}}{\text{chance variation}}
\]

The test statistic depends on

\begin{itemize}
  \item The distance between the estimate and the null value
  \item The population variability $\sigma$
  \item The sample size $n$
\end{itemize}

\vspace{0.2cm}
A very large sample can make a small effect statistically significant. A small or highly variable sample may fail to detect an effect that matters in practice.
\end{frame}

\begin{frame}{The Same Effect at Different Sample Sizes}
\centering
\lectureimage[0.88\linewidth]{fig_ch15_significance_sample_size.png}
\end{frame}

% SLIDE 22
\begin{frame}{Statistical Significance and Practical Importance}
Statistical significance addresses whether the data are difficult to explain under the null model.

\vspace{0.2cm}
Practical importance addresses whether the estimated effect is large enough to matter in the scientific, clinical or policy setting.

\begin{center}
\begin{tabularx}{0.92\textwidth}{>{\RaggedRight\arraybackslash}p{0.28\textwidth} X}
\toprule
\textbf{Question} & \textbf{Useful information} \\
\midrule
Is there evidence of an effect? & \textit{P}-value and model conditions \\
How large might the effect be? & Effect estimate and confidence interval \\
Does the effect matter? & Subject-matter thresholds, costs, benefits and consequences \\
\bottomrule
\end{tabularx}
\end{center}

A complete analysis reports the size and uncertainty of the effect, not only whether a threshold was crossed.
\end{frame}

% SLIDE 23
\begin{frame}{What a $P$-Value Does Not Mean}
A \textit{P}-value is calculated under the null hypothesis and the assumptions of the statistical model.

\vspace{0.2cm}
It is not

\begin{itemize}
  \item The probability that $H_0$ is true
  \item The probability that the observed result occurred only by chance
  \item The size or importance of the observed effect
  \item Evidence that can be interpreted without checking the study design and model conditions
\end{itemize}

\vspace{0.2cm}
A large \textit{P}-value means that the data remain compatible with the null model. It does not establish that the null hypothesis is correct.
\end{frame}

% SLIDE 24
\begin{frame}{Multiple Analyses and Selective Reporting}
Each test conducted at level $\alpha$ has a chance of producing a Type I error when its null hypothesis is true.

\vspace{0.2cm}
Repeatedly testing many outcomes, subgroups or models increases the chance that at least one result will appear statistically significant by chance.

\begin{block}{Illustration}
If $20$ independent tests are conducted at $\alpha=0.05$ and all null hypotheses are true, then
\[
P(\text{at least one rejection})
=1-(0.95)^{20}
\approx0.642
\]
\end{block}

Searching until a small \textit{P}-value appears is often called \textit{P}-hacking or data dredging. Reporting only significant findings also creates publication bias.
\end{frame}

\begin{frame}{False Positives Across Multiple Tests}
\centering
\lectureimage[0.88\linewidth]{fig_ch15_multiple_testing.png}
\end{frame}

% SLIDE 25
\begin{frame}{Inference Within the Scientific Process}
Statistical inference supports scientific reasoning when it forms part of a transparent process.

\begin{enumerate}
  \item State the research question and analysis plan before examining the results
  \item Use a design that connects the data to the intended population or causal question
  \item Examine the data and assess the conditions for the procedure
  \item Report estimates, uncertainty and exact \textit{P}-values with appropriate context
  \item Replicate important findings with new data when possible
\end{enumerate}

\vspace{0.2cm}
Independent replication reduces the chance that a conclusion depends on an unusual sample, an unnoticed design problem or a flexible analysis choice.
\end{frame}

% =============================================================
% THEORY: PLANNING STUDIES
% =============================================================

\section{Planning Studies}

% SLIDE 26
\begin{frame}{Planning a Confidence Interval}
For a one-sample $z$ confidence interval,

\[
m=z^*\frac{\sigma}{\sqrt{n}}
\]

Solving for the sample size required to achieve a desired margin of error gives

\[
n
=
\left(\frac{z^*\sigma}{m}\right)^2
\]

\vspace{0.2cm}
The calculated value must be rounded \textbf{up} so that the achieved margin of error is no larger than the target.

\vspace{0.2cm}
The calculation requires a plausible value of $\sigma$, often obtained from previous studies or a pilot study.
\end{frame}

% SLIDE 27
\begin{frame}{Population Size and the Required Sample}
The usual sample-size formula depends on

\begin{itemize}
  \item The chosen confidence level through $z^*$
  \item The expected population variability through $\sigma$
  \item The desired precision through $m$
\end{itemize}

\vspace{0.2cm}
The population size does not appear in the formula when the population is much larger than the sample.

\begin{block}{Interpretation}
A sample of a particular size can estimate a mean with the same planned margin of error in two very large populations, provided that the sampling design and population variability are comparable.
\end{block}

Special corrections may be used when the sample is a substantial fraction of a finite population.
\end{frame}

% SLIDE 28
\begin{frame}{Planning a Hypothesis Test}
Planning a hypothesis test requires three choices.

\begin{itemize}
  \item \textbf{Significance level $\alpha$:} the tolerated probability of a Type I error when $H_0$ is true
  \item \textbf{Effect size:} the smallest departure from the null value that would matter in context
  \item \textbf{Power:} the desired probability that the test will detect that effect
\end{itemize}

\vspace{0.25cm}
These choices should be made before data collection. They connect the statistical design to the scientific consequences of missing or incorrectly declaring an effect.
\end{frame}

% SLIDE 29
\begin{frame}{Power and the Probability of a Type II Error}
The \textbf{power} of a test against a specified alternative value is the probability that the test rejects $H_0$ when that alternative value is true.

\[
\text{Power}=P(\text{reject }H_0\mid\text{specified alternative is true})
\]

If $\beta$ is the probability of a Type II error for that alternative, then

\[
\text{Power}=1-\beta
\]

\vspace{0.2cm}
Power generally increases with a larger sample, a larger effect, lower population variability or a less restrictive significance level. Two-sided tests generally require more observations than one-sided tests with the same target power.
\end{frame}

\begin{frame}{Effect Size and Required Sample Size}
\centering
\lectureimage[0.88\linewidth]{fig_ch15_power_sample_size.png}
\end{frame}

% =============================================================
% EXAMPLES
% =============================================================

\section{Examples}

% SLIDE 30
\begin{frame}{Example 15.1: Two Convenience Samples}
Two researchers recruit students enrolled in introductory courses at one university.

\begin{block}{Neurobiology Study}
A neurobiologist studies how people with normal vision respond to an optical illusion. Basic visual mechanisms may be similar across many groups, so subject-matter knowledge may support a cautiously broader interpretation.
\end{block}

\begin{block}{Sociology Study}
A sociologist studies attitudes toward using human participants in research. Age, education, field of study and campus culture may all influence the response, so one sociology class is unlikely to represent adults or university students generally.
\end{block}

Neither class is an actual random sample. The plausibility and scope of generalization differ because the research questions differ.
\end{frame}

% SLIDE 31
\begin{frame}{Example 15.2: Evaluating the Source of the Data}
Three Chapter 14 datasets illustrate different levels of support for inference.

\begin{center}
\small
\begin{tabularx}{0.96\textwidth}{>{\RaggedRight\arraybackslash}p{0.22\textwidth} >{\RaggedRight\arraybackslash}p{0.28\textwidth} X}
\toprule
\textbf{Study} & \textbf{Data source} & \textbf{Practical assessment} \\
\midrule
NHANES heights & Complex multistage probability sample & Broadly representative, but the formal analysis should account for the design \\
School-district IQ scores & Random sample of girls in the defined district & Supports inference to the stated district population \\
Blood phosphorus & Random records from one Ohio laboratory & Supports a narrow laboratory-related population, not all adults aged $75$ to $79$ \\
\bottomrule
\end{tabularx}
\end{center}

The sampling frame and measurement setting determine the population to which the result can reasonably apply.
\end{frame}

% SLIDE 32
\begin{frame}{Example 15.3: Confidence Level and Margin of Error}
For the IQ study,

\[
\bar{x}=105.84
\qquad
\sigma=15
\qquad
n=31
\]

The margin of error changes with the confidence level.

\begin{center}
\begin{tabular}{cccc}
\toprule
\textbf{Confidence} & \textbf{$z^*$} & \textbf{Margin} & \textbf{Interval} \\
\midrule
$90\%$ & $1.645$ & $4.43$ & $(101.41,\ 110.27)$ \\
$95\%$ & $1.960$ & $5.28$ & $(100.56,\ 111.12)$ \\
$99\%$ & $2.576$ & $6.94$ & $(98.90,\ 112.78)$ \\
\bottomrule
\end{tabular}
\end{center}

All three intervals have the same centre. Higher confidence produces a wider interval because it uses a larger critical value.
\end{frame}

\begin{frame}{Comparing Confidence Levels}
\centering
\lectureimage[0.88\linewidth]{fig_ch15_confidence_levels.png}
\end{frame}

% SLIDE 33
\begin{frame}{Example 15.3: Increasing the Sample Size}
At $95\%$ confidence, the original margin of error is

\[
m_{31}
=1.960\frac{15}{\sqrt{31}}
\approx5.28
\]

If the sample size increases from $31$ to $62$ while $sigma$ remains $15$, then

\[
m_{62}
=1.960\frac{15}{\sqrt{62}}
\approx3.73
\]

\[
\frac{m_{62}}{m_{31}}
=\sqrt{\frac{31}{62}}
=\frac{1}{\sqrt{2}}
\approx0.707
\]

Doubling the sample size reduces the margin by about $29.3\%$. A new sample would usually also produce a different value of $\bar{x}$, so the centre may change.
\end{frame}

% SLIDE 34
\begin{frame}{Example 15.4: A Large Social-Media Experiment}
A randomized experiment altered the amount of emotional content shown in the News Feeds of $689{,}003$ Facebook users for one week.

\vspace{0.2cm}
The reported tests produced very small \textit{P}-values:

\[
z=-4.78
\qquad\text{and}\qquad
z=-7.219
\qquad
P<0.001
\]

The estimated change was approximately one fewer emotional word per $1000$ words written during the following week.

\begin{block}{Interpretation}
The large sample made the analysis sensitive to a very small average change. The small \textit{P}-values provide strong evidence against a zero-effect model, but they do not show that the effect was large.
\end{block}
\end{frame}

% SLIDE 35
\begin{frame}{Example 15.5: Assessing Practical Importance}
The Facebook experiment raises questions that a \textit{P}-value cannot answer.

\begin{itemize}
  \item Does using fewer emotional words reflect a meaningful change in emotion?
  \item Could users simply repeat words that recently appeared in their feeds?
  \item Would the experimental manipulation resemble ordinary use of the platform?
  \item Could a very small individual effect matter when applied across a large population?
\end{itemize}

\vspace{0.2cm}
Evaluating practical importance requires the estimated effect, its uncertainty, the quality of the measurement and knowledge of the setting.

\begin{block}{Conclusion}
Statistical significance answers a narrow probability question. Subject-matter reasoning determines whether the finding is meaningful.
\end{block}
\end{frame}

% SLIDE 36
\begin{frame}{Example 15.6: Planning a Body-Temperature Study}
Researchers want to estimate the mean body temperature of healthy adults within $0.05^\circ$F at $95\%$ confidence. Previous research suggests

\[
\sigma=0.6^\circ\mathrm{F}
\qquad\text{and}\qquad
z^*=1.960
\]

The required sample size is

\[
n
=
\left(\frac{1.960(0.6)}{0.05}\right)^2
=553.1904
\]

The study must recruit at least

\[
\boxed{n=554}
\]

Rounding down to $553$ would produce a margin of error slightly larger than the target.
\end{frame}

% SLIDE 37
\begin{frame}{Example 15.7: Planning a Sweetness Study}
A quality-control engineer studies the mean loss of sweetness after cola is stored. Each trained taster gives a before-and-after rating.

\[
H_0:\mu=0
\qquad
H_a:\mu>0
\]

The planning values are

\begin{itemize}
  \item Population standard deviation $\sigma=1$ point
  \item Significance level $\alpha=0.05$
  \item Important mean loss $\mu=0.8$ point
  \item Desired power $0.90$
\end{itemize}

The initial design uses $n=10$ tasters. The next step is to determine whether this sample provides the target power against $\mu=0.8$.
\end{frame}

% SLIDE 38
\begin{frame}{Example 15.8: Calculating Power for $n=10$}
For an upper-tail test with $\alpha=0.05$, the critical sample mean under $H_0$ is

\[
\bar{x}_{\mathrm{crit}}
=0+1.645\left(\frac{1}{\sqrt{10}}\right)
\approx0.520
\]

When the true mean loss is $\mu=0.8$, the power is

\[
P(\bar{X}\ge0.520\mid\mu=0.8)
=
P\left(
Z\ge
\frac{0.520-0.8}{1/\sqrt{10}}
\right)
\]

\[
=P(Z\ge-0.886)
\approx0.812
\]

The test detects a true mean loss of $0.8$ point about $81.2\%$ of the time. Ten tasters do not provide the desired $90\%$ power.
\end{frame}

\begin{frame}{Null and Alternative Sampling Distributions}
\centering
\lectureimage[0.88\linewidth]{fig_ch15_power_curves.png}
\end{frame}

% SLIDE 39
\begin{frame}{Example 15.9: Error Probabilities in the Sweetness Study}
The test uses $\alpha=0.05$ and has power $0.812$ against the specific alternative $\mu=0.8$.

\begin{block}{Type I Error}
\[
P(\text{reject }H_0\mid\mu=0)
=\alpha
=0.05
\]
The study declares a sweetness loss even though the population mean loss is $0$.
\end{block}

\begin{block}{Type II Error at $\mu=0.8$}
\[
\beta
=1-\text{power}
=1-0.812
=0.188
\]
The study fails to detect a mean loss of $0.8$ point when that loss is real.
\end{block}
\end{frame}

% SLIDE 40
\begin{frame}{Chapter 15 Summary}
Reliable inference requires more than substituting numbers into a formula.

\begin{itemize}
  \item The data-collection process determines the scope of generalization and causal interpretation
  \item Confidence intervals become wider with higher confidence or greater variability, and narrower with larger samples
  \item The margin of error measures random sampling variation rather than every possible source of error
  \item Statistical significance depends on effect size, variability and sample size
  \item A \textit{P}-value does not measure the probability that $H_0$ is true or the practical importance of an effect
  \item Study planning connects sample size with precision, significance level, meaningful effect size and power
\end{itemize}

\vspace{0.2cm}
Clear conclusions report the estimate, uncertainty, study limitations and practical meaning together.
\end{frame}

\end{document}
